\chapter*{Abstract}
\addcontentsline{toc}{chapter}{Abstract}

Automation in football decisions attempts to reduce subjective error to protect the integrity of the game's most crucial moments. However, the infrastructure required to support offside decisions remains unavailable outside the elite tier of the sport. This final year project investigates the feasibility of automated offside detection from broadcast video alone, through the design and evaluation of a modular computer vision pipeline.

The pipeline operates on broadcast clips drawn from the publicly available SoccerNet dataset, covering thousands of games across the top five European Leagues and even the UEFA Champions League up until the 2016/2017 season. For each offside clip in each downloaded game, the keyframe corresponding to the moment of the pass is manually selected via a browser-based annotation interface. Players are then detected using a YOLO-based object detector, and team identification is done through clustering kit colours in Lab colour space using K-Means. A homography between the broadcast image plane and a real-world pitch coordinate system is estimated using RANSAC over manually picked relevant pitch landmarks, allowing every player's position to be projected onto a metric pitch. The offside condition is then evaluated geometrically by comparing the attacking player against the deepest outfield defender along the direction of attack.

One hundred and fifty offside events were processed and compared against the annotations provided by SoccerNet. Of these, 128 verdicts (85.3\%) agreed with the dataset annotations. Since the annotations reflect the match official's judgement rather than a geometric ground truth, disagreements may originate from system error or from official error; the two cannot be distinguished without independent measurement. Agreement rose to 91.3\% for the 92 events where the measured margin was above 0.5\,m, the threshold below which both human judgement and the system's geometric precision become questionable, demonstrating that broadcast-only offside detection is feasible on representative footage. The principal sources of system error include occasional team misidentification and imprecision in the estimated homography.

The major computational expense of the automation is the two YOLO11x-seg inference passes per event, which take around 90\,ms on commodity GPU hardware. The main remaining obstacles to fully automated deployment are the manual keyframe selection and landmark picking stages, which are identified as the primary directions for future work.